\documentclass[11pt,a4paper]{article}
\usepackage[T1]{fontenc}
\usepackage{isabelle,isabellesym}

% further packages required for unusual symbols (see also
% isabellesym.sty), use only when needed


\usepackage{amssymb,amsmath}
  %for \<leadsto>, \<box>, \<diamond>, \<sqsupset>, \<mho>, \<Join>,
  %\<lhd>, \<lesssim>, \<greatersim>, \<lessapprox>, \<greaterapprox>,
  %\<triangleq>, \<yen>, \<lozenge>

%\usepackage{eurosym}
  %for \<euro>

\usepackage[only,bigsqcap]{stmaryrd}
%for \<Sqinter>

%\usepackage{eufrak}
  %for \<AA> ... \<ZZ>, \<aa> ... \<zz> (also included in amssymb)

%\usepackage{textcomp}
  %for \<onequarter>, \<onehalf>, \<threequarters>, \<degree>, \<cent>,
  %\<currency>

% this should be the last package used
\usepackage{pdfsetup}

% urls in roman style, theory text in math-similar italics
\urlstyle{rm}
\isabellestyle{it}


% for uniform font size
%\renewcommand{\isastyle}{\isastyleminor}


\begin{document}

\title{Cauchy's Functional Equation}
\author{Michikazu Hirata}
\maketitle
\begin{abstract}
  Cauchy's functional equation is the equation
  \[h(x + y) = h(x) + h(y).\]
  This entry formalizes properties of the solutions
  $h:\mathbb{R}\to\mathbb{R}$ of Cauchy's functional equation, including:
  \begin{itemize}
  \item A function $h$ satisfying Cauchy's functional equation has the form
        $h(x) = c x$ for some $c$ if one of the following conditions holds:
    \begin{itemize}
    \item $h$ is continuous at some point.
    \item $h$ is continuous.
    \item $h$ is bounded on some interval.
    \item $h$ is (anti)monotone on some interval.
    \item $h$ is measurable.
    \end{itemize}
  \item There are non-linear solutions to Cauchy's functional equation.
  \item The set of non-linear solutions has cardinality $|\mathbb{R}^{\mathbb{R}}|$.
  \item The graph of a non-linear solution $\{(x,h(x))\mid x\in\mathbb{R}\}$ is
  dense in $\mathbb{R}\times\mathbb{R}$.
  \end{itemize}
\end{abstract}

\tableofcontents

% sane default for proof documents
\parindent 0pt\parskip 0.5ex

% generated text of all theories
\input{session}

% optional bibliography
\bibliographystyle{abbrv}
\bibliography{root}

\end{document}

%%% Local Variables:
%%% mode: latex
%%% TeX-master: t
%%% End:
